\section{מבוא ושאלות המחקר}

סדרות זמן פיננסיות אינן מתנהגות בהכרח כתהליך יחיד ויציב. תקופות שקטות מתחלפות בתקופות תנודתיות, התפלגות התשואות משתנה, והקשר בין נכסים יכול להתחזק או להיחלש. מודלים מסוג regime switching מתארים מצב זה באמצעות משתנה חבוי שמייצג את הסביבה הסטטיסטית שבה השוק נמצא \cite{hamilton1989new}. מודל מרקוב חבוי נותן לכך מסגרת הסתברותית פשוטה: state סמוי מתפתח לפי שרשרת מרקוב, והתצפיות נוצרות מהתפלגות שתלויה ב־state \cite{rabiner1989tutorial,murphy2012machine}.

הפרויקט התחיל משאלה יחסית ישירה: האם \tech{Gaussian HMM} הלומד מצבי שוק סמויים יכול לשפר חיזוי של כיוון התשואה ביום הבא, וכיצד מספר המצבים \(K\) משפיע על התוצאה. רעיון דומה מופיע בעבודות המיישמות HMM לנתוני מניות לצורך forecasting \cite{catello2023hmm}. עם זאת, לאחר הריצה הראשונית התברר שהשאלה החיזויית לבדה צרה מדי: גם אם DPA אינו גבוה, ייתכן שהמודל עדיין מזהה מבנה סטטיסטי חשוב שאינו נגיש ל־classifier רגיל.

לכן בגרסה הסופית הוגדרו שלוש שאלות מחקר:

\begin{enumerate}
    \item \textbf{Forecasting:} האם HMM משפר חיזוי next-day direction ביחס ל־Baselines פשוטים ולמודלי Supervised ML המקבלים את אותם Features?
    \item \textbf{Representation:} האם ה־Hidden States נבדלים באופן עקבי ב־return, volatility, daily range, volume behavior, drawdown, occupancy ו־persistence?
    \item \textbf{Regime information:} האם state identity וה־posterior uncertainty מספקים מידע נוסף על מעברי משטר, VIX ומבנה הקורלציה בין נכסים?
\end{enumerate}

ההבחנה בין forecasting לבין representation היא מרכזית. מודל יכול להיות חלש בחיזוי נקודתי ועדיין להיות שימושי כאומד של מצב חבוי. לכן העבודה אינה מחפשת רק ``מודל מנצח'', אלא בודקת איזה סוג מידע כל משפחת מודלים מספקת ומה ניתן להסיק ממנו באופן זהיר.
